% ============================================================
% 片段：s04b  对应原文 PDF 第 32--38 页
% 内容：§4.4 高亏格振幅的多项式结构；§4.5 使用多项式结构（4.5.1--4.5.3）；致谢
% 译者：t-s04b   审校状态：初稿
% 说明：原文 4.5.2、4.5.3 中的不编号斜体小标题（Genus 1、Higher genus boundary
%       conditions 等）在原文中为 \paragraph 级斜体标题，此处以斜体独立行复现。
% ============================================================

% ==== 原文 PDF 第 32 页（§4.4 起） ====

\subsection{高亏格振幅的多项式结构}
\label{sec:4.4}

\begin{keytakeaways}
  \begin{itemize}[leftmargin=1.4em,itemsep=2pt,topsep=2pt]
    \item \textbf{物理图像}：拓扑弦的高亏格自由能可解释为物理弦理论的某些振幅；完整的拓扑弦配分函数
          还被猜想编码了 4d BPS 态的信息 \cite{ref73}，因此自由能理应带有靶空间对偶群的自守形式结构。
    \item \textbf{数学工具}：特殊几何，具体说是出现在全纯反常方程右端的联络多重导数。Yamaguchi 与
          Yau 用 $A_p=G^{z\bar z}(z\partial_z)^pG_{z\bar z}$、$B_p=e^{K}(z\partial_z)^pe^{-K}$ 与全纯
          生成元 $X\sim1/\Delta$ 组成生成元环，并把它约化为 $A_1,B_1,B_2,B_3$ 与 $X$。
    \item \textbf{关键结论}：复结构形变空间为一维时，高亏格振幅是有限个生成元中次数为 $3g-3+n$ 的
          多项式 \cite{ref21}；任意维模空间的情形由生成元 $K_i,S^{ij},S^i,S$（分次 $1,1,2,3$）实现，式 (4.29)
          的换元使振幅不再依赖 $\tilde K_i$ \cite{ref22}。
    \item \textbf{技术核心}：反常方程在生成元的求导下封闭（式 (4.32)），故每一步递推只是把低亏格的
          已知量代数地组合起来；剩下的全纯模糊性由第~\ref{sec:4.5.2}~节的边界条件逐阶固定。
  \end{itemize}
\end{keytakeaways}

% [Tutor Note: 物理图像] 多项式结构的物理根源：非全纯部分由低亏格骨架唯一决定
\begin{yizhu}[物理图像]
“多项式结构”之所以可能，是因为反常方程右端只含有限几类量：汤川耦合 $C_{ijk}$、凯勒势 $K$ 及其
导数、以及联络 $\Gamma^l_{ij}$ 的若干组合。递推时，$\bar\partial_{\bar\imath}\mathcal{F}^{(g)}$
被这些“原子”乘上低亏格振幅的组合表示出来，因此 $\mathcal{F}^{(g)}$ 的非全纯依赖永远是有限个生成元
的多项式；新的物理信息只藏在全纯积分常数里。这正是“高亏格振幅由低亏格骨架决定”的精确表述，也解释
了为何振幅在对偶群下的模性足以用来确定反常方程的解 \cite{ref19}。
\end{yizhu}

拓扑弦的高亏格自由能可以解释为给出物理弦理论的某些振幅。\footnote{综述见文献 \cite{ref72}。} 完整的
拓扑弦配分函数被猜想也编码了 4d BPS 态的信息 \cite{ref73}。因此很自然地期望，%
% ==== 原文 PDF 第 33 页 ====
拓扑弦的自由能由靶空间对偶群的自守形式（automorphic forms）刻画。文献 \cite{ref19} 利用拓扑弦振幅的模性
（modularity）来确定反常方程的解。每当模群是 $\mathrm{SL}(2,\mathbb{Z})$ 或其子群时，振幅的模性
表现得最为明显。椭圆曲线的 Gromov--Witten 不变量的高亏格生成函数曾被写成多项式 \cite{ref74,ref75}，其多项式
生成元是近全纯模形式（almost holomorphic modular forms）环的元素 $E_2$、$E_4$ 与 $E_6$ \cite{ref76}。只要
$\mathrm{SL}(2,\mathbb{Z})$ 是模群的子群，这类生成元的多项式也会出现，例如见文献
\cite{ref77,ref78,ref79,ref80,ref28,ref81}。拓扑弦与近全纯模形式的关系在文献 \cite{ref23,ref82,ref24} 中得到进一步探讨。然而，
对于大多数 CY 流形的模空间及其单值群（monodromy group），相应的自守形式并不知道。结果表明，仅用
特殊几何就足以证明高亏格拓扑弦振幅的多项式结构 \cite{ref21,ref22}，下面将对此加以综述。

Yamaguchi 与 Yau 在 \cite{ref21} 中证明：对于镜像五次超曲面以及复结构形变空间为一维的相关几何，拓扑弦的
高亏格振幅可以表示为有限个生成元中次数为 $3g-3+n$ 的多项式，其中 $n$ 指插入算符的个数。作为多项式
的生成元，他们使用了特殊几何中的某些元素，也就是出现在全纯反常方程右端的那些量——联络的多重导数。
文献 \cite{ref21} 的非全纯生成元由
\[
  A_p = G^{z\bar z}(z\partial_z)^p G_{z\bar z},
  \qquad
  B_p = e^{K}(z\partial_z)^p e^{-K},
  \qquad p=1,2,3,\ldots
\]
给出。此外还使用了一个全纯生成元 $X\sim 1/\Delta$，其中 $\Delta$ 为判别式。下一步要证明生成元之间的
关系，使得无穷多个非全纯生成元可以约化为 $A_1$、$B_1$、$B_2$ 与 $B_3$。进一步，由全纯反常方程的
分析可以证明 \cite{ref21}：拓扑弦振幅中只出现生成元的特殊组合，因而非全纯生成元的数目又减少一个。

% [Tutor Note: 含糊点] “生成元之间的关系”与“特殊组合”两句跳过了 \cite{ref21} 的核心技术步骤
\begin{yizhu}[含糊点]
原文只说“证明生成元之间的关系”而把无穷多个 $A_p,B_p$ 约化为 $A_1,B_1,B_2,B_3$，又说“只出现特殊
组合”使数目再减一，但未写出这些关系式。就 \cite{ref21} 的构造而言，约化依赖于 $A_p,B_p$ 之间由特殊几何导出
的微分关系，以及反常方程中这些量只以特定组合出现这一事实；具体关系式原文未给。核对时应直接查文献
\cite{ref21} 的生成元关系式，不要试图由本讲义的叙述推出。
\end{yizhu}

将这一构造推广到模空间维数任意的 CY 流形，是在文献 \cite{ref22} 中给出的。\footnote{这一问题由 P. Mayr
在一份未发表的手稿中独立解决。} 下面将综述这一构造，以及为在递推的每一步求解全纯模糊性所需的边界
条件。

文献 \cite{ref22} 证明：关联函数 $\mathcal{F}^{(g)}_{i_1\cdots i_n}$ 是生成元 $K_i$、$S^{ij}$、$S^i$、$S$
中次数为 $3g-3+n$ 的多项式，这些生成元的分次（grading）分别取为 $1,1,2,3$。文献 \cite{ref22} 进一步证明，
通过如下生成元变换
\begin{equation}\tag{4.29}
\begin{aligned}
  \tilde S^{ij} &= S^{ij}, \\
  \tilde S^i &= S^i - S^{ij}K_j, \\
  % ==== 原文 PDF 第 34 页 ====
  \tilde S &= S - S^iK_i + \frac{1}{2}S^{ij}K_iK_j, \\
  \tilde K_i &= K_i,
\end{aligned}
\end{equation}
$\mathcal{F}^{(g)}$ 不依赖于 $\tilde K_i$，即 $\partial\mathcal{F}^{(g)}/\partial\tilde K_i=0$。
下文将去掉修改后生成元上的波浪号。

% [Tutor Note: 含糊点] (4.29) 中的换元与 (4.34) 中 $\partial/\partial K_j$ 项的基准问题
\begin{yizhu}[含糊点]
式 (4.29) 之后紧接着断言 $\partial\mathcal{F}^{(g)}/\partial\tilde K_i=0$，而式 (4.34) 右端却保留
$G_{\bar\imath j}\partial\mathcal{F}^{(g)}/\partial K_j$。这一点原文没有展开：它并未交代 (4.34)
写在未加波浪号还是已加波浪号的那一套生成元下，也未说明该断言的证明。阅读时只要记住 (4.29) 是一次
保持分次的生成元变换（$\tilde S$ 的表达式相当于把 $S$ 沿 $S^i$ 方向作平移 $K$ 后“配方”），
$K$ 的依赖被平移进 $\tilde S^i$ 与 $\tilde S$ 之中；(4.34) 则按原文形式照录，未作改写。核对时请以
文献 \cite{ref22} 的原始推导为准。
\end{yizhu}

证明的出发点是：把最先出现的非零关联函数用这些生成元表示。在亏格零，这些就是全纯三点耦合
$\mathcal{F}^{(0)}_{ijk}=C_{ijk}$。全纯反常方程 (4.24) 可以利用 (4.27) 积分，得到
\begin{equation}\tag{4.30}
\mathcal{F}^{(1)}_i
= \frac{1}{2}C_{ijk}S^{jk} + \left(1-\frac{\chi}{24}\right)K_i + f^{(1)}_i,
\end{equation}
其中 $f^{(1)}_i$ 为模糊性。由这一表达式可以看出，关联函数的非全纯性只来自生成元。此外，特殊几何关系
(4.11) 可以积分为
\begin{equation}\tag{4.31}
\Gamma^l_{ij} = \delta^l_i K_j + \delta^l_j K_i - C_{ijk}S^{kl} + s^l_{ij},
\end{equation}
其中 $s^l_{ij}$ 表示未被特殊几何关系固定的全纯函数；由此可以导出下列方程，它们表明承载非全纯性的
生成元在求导下封闭 \cite{ref22}。\footnote{这些方程是对式 (4.29) 中带波浪号的生成元写出的，可由文献 \cite{ref22}
中的方程直接得到。}
\begin{equation}\tag{4.32}
\begin{aligned}
  \partial_i S^{jk} &= C_{imn}S^{mj}S^{nk} + \delta^j_i S^k + \delta^k_i S^j
      - s^j_{im}S^{mk} - s^k_{im}S^{mj} + h_i^{jk}, \\
  \partial_i S^j &= C_{imn}S^{mj}S^n + 2\delta^j_i S
      - s^j_{im}S^m - h_{ik}S^{kj} + h^j_i, \\
  \partial_i S &= \frac{1}{2}C_{imn}S^m S^n - h_{ij}S^j + h_i, \\
  \partial_i K_j &= K_iK_j - C_{ijn}S^{mn}K_m + s^m_{ij}K_m - C_{ijk}S^k + h_{ij},
\end{aligned}
\end{equation}
其中 $h_i^{jk}$、$h_i^j$、$h_i$ 与 $h_{ij}$ 表示全纯函数。所有这些函数与式 (4.31) 中的函数并非彼此
独立。文献 \cite{ref26} 证明（另见 \cite{ref25}）：在此环中选取全纯函数的自由度归结为函数 $\mathcal{E}^{ij}$、
$\mathcal{E}^j$ 与 $\mathcal{E}$，它们可以加到多项式生成元上
\begin{equation}\tag{4.33}
\begin{aligned}
  \widehat{S}^{ij} &= S^{ij} + \mathcal{E}^{ij}, \\
  \widehat{S}^j &= S^j + \mathcal{E}^j, \\
  \widehat{S} &= S + \mathcal{E}.
\end{aligned}
\end{equation}
所有全纯量都相应地改变。

% [Tutor Note: 易错点] $S^{ij}$（生成元）与 $s^l_{ij}$（全纯系数）是两类完全不同的对象
\begin{yizhu}[易错点]
式 (4.31)--(4.33) 中大小写极易混淆：$S^{ij}$ 是非全纯的传播子型生成元，而 $s^l_{ij}$ 是用于重写联络
的全纯函数；同理 $\Gamma^l_{ij}$ 含非全纯部分，$\mathcal{E}^{ij}$、$\mathcal{E}^j$、$\mathcal{E}$ 则
纯全纯。核对时按“是否含反全纯指标”区分：$s$、$h$、$\mathcal{E}$ 都是全纯的待定函数，它们的选取对应
(4.33) 的自由度；而 $S^{ij},S^i,S$ 一旦被定义为生成元，其反全纯导数就由 (4.32) 唯一给出。
\end{yizhu}

% [Tutor Note: 易错点] 生成元的分次与多项式次数 $3g-3+n$ 的核算
\begin{yizhu}[易错点]
原文说生成元 $K_i,S^{ij},S^i,S$ 的分次（grading）依次为 $1,1,2,3$，注意 $S^{ij}$ 的分次也是 $1$，
并非 $2$。可用式 (4.30) 自检：$g=1$、$n=1$ 时总次数应为 $3g-3+n=1$，而式中 $C_{ijk}S^{jk}$ 与
$K_i$ 都恰为一次。若误把 $S^{ij}$ 记为二次，则后续所有“次数为 $3g-3+n$”的计数（以及 (4.37)、
(4.38) 的量纲）都会出错。
\end{yizhu}

现在，拓扑弦振幅满足的全纯反常方程中，$\bar\partial_{\bar\imath}$ 导数被替换为对多项式生成元的
导数 \cite{ref28,ref22}。
\begin{equation}\tag{4.34}
\bar\partial_{\bar\imath}\mathcal{F}^{(g)}
= \bar{C}^{jk}_{\bar\imath}\Bigl\{
  \frac{\partial\mathcal{F}^{(g)}}{\partial S^{jk}}
  - \frac{1}{2}\frac{\partial\mathcal{F}^{(g)}}{\partial S^k}K_j
  - \frac{1}{2}\frac{\partial\mathcal{F}^{(g)}}{\partial S^j}K_k
  + \frac{1}{2}\frac{\partial\mathcal{F}^{(g)}}{\partial S}K_jK_k \Bigr\}
+ G_{\bar\imath j}\frac{\partial\mathcal{F}^{(g)}}{\partial K_j}
\end{equation}
% ==== 原文 PDF 第 35 页 ====
\begin{equation}\tag{4.35}
\begin{aligned}
&= \frac{1}{2}\bar{C}^{jk}_{\bar\imath}\Bigl\{
  \Bigl(\partial_j\partial_k + (C_{jkm}S^{ml} - s^l_{jk})\partial_l
  + (4-2g)(h_{jk} - C_{jkl}S^l) \\
&\qquad + (3-2g)K_j\partial_k + (3-2g)K_k\partial_j
  + (4-2g)(3-2g)K_kK_j\Bigr)\mathcal{F}^{(g-1)} \\
&\quad + \sum_{h=1}^{g-1}\partial_j\mathcal{F}^{(g-h)}\partial_k\mathcal{F}^{(h)}
  + (2-2h)K_j\mathcal{F}^{(h)}\partial_k\mathcal{F}^{(g-h)}
  + (2-2g+2h)K_k\mathcal{F}^{(g-h)}\partial_j\mathcal{F}^{(h)} \\
&\qquad + K_jK_k(2-2h)(2-2g+2h)\mathcal{F}^{(h)}\mathcal{F}^{(g-h)}
  \Bigr\}.
\end{aligned}
\end{equation}

\subsection{使用多项式结构}
\label{sec:4.5}

\begin{keytakeaways}
  \begin{itemize}[leftmargin=1.4em,itemsep=2pt,topsep=2pt]
    \item \textbf{物理图像}：多项式结构把“解无穷多个亏格的反常方程”化归为“有限个待定系数的代数问题”：
          每个亏格的 $\mathcal{F}^{(g)}$ 写成生成元的多项式，系数是模参数的函数，再由边界条件逐步确定。
    \item \textbf{数学工具}：取局部坐标 $z_*$ 使 $C_{*ij}$ 可逆，由此从特殊几何关系反解出 $S^{ij}$
          （式 (4.36)）；再用闭合方程 (4.32) 逐次构造 $S^i$、$S$（式 (4.37)、(4.38)）。边界条件有三类：
          大复结构极限的常映射贡献 (4.42)、类锥形流形处的奇异结构 (4.43)，以及全纯模糊性的局部 ansatz
          (4.44)——后者的有限个系数由前两类条件固定。
    \item \textbf{关键结论}：生成元与 $s,h$ 的全纯规范选定后，高亏格振幅的全部模糊性化为
          $\Delta^{-(2g-2)}$ 乘一个低次多项式 $p(z)$ 的有限个系数；这些系数被 (4.42)、(4.43) 等边界
          条件线性地固定。
    \item \textbf{例}：五次超曲面在 $g=1$ 由 (4.46) 的 $n,\chi,s,r$ 完全确定（式 (4.47)）；$g=2$ 的
          (4.48) 给出多项式结构的一次显式检验。
  \end{itemize}
\end{keytakeaways}

% [Tutor Note: 物理图像] 4.5 节三步走：造生成元、给边界条件、在五次超曲面上验证
\begin{yizhu}[物理图像]
本节把上一小节的结论变成可执行的算法，分三步：(1) 构造生成元，即在模空间的局部坐标下把
$S^{ij},S^i,S$ 具体解出来（4.5.1）；(2) 给出边界条件，用大复结构极限、类锥形流形轨迹与全纯模糊性的
ansatz 固定每一步的积分常数（4.5.2）；(3) 在五次超曲面上显式执行到 $g=2$（4.5.3）。之所以能做第三
步，是因为多项式结构保证每个 $g$ 只有有限个未知系数，而边界条件恰好给出同样多的线性方程。
\end{yizhu}

\subsubsection{构造生成元}
\label{sec:4.5.1}

\begin{keytakeaways}
  \begin{itemize}[leftmargin=1.4em,itemsep=2pt,topsep=2pt]
    \item \textbf{物理图像}：$S^{ij}$、$S^i$、$S$ 就是反常方程右端那些非全纯数据的“积分”，在费曼图
          图像中充当传播子；构造生成元即选定一组传播子与全纯规范。
    \item \textbf{数学工具}：取局部坐标 $z_*$ 使 $C_{*ij}$ 可逆，从特殊几何关系 (4.31) 反解出 $S^{ij}$
          （式 (4.36)），再用闭合方程 (4.32) 逐次解出 $S^i$、$S$（式 (4.37)、(4.38)）。
    \item \textbf{关键结论}：一旦全纯函数 $s$、$h$ 按 (4.33) 的自由度选定，全部生成元就确定下来，
          且所有全纯函数都化为局部坐标 $z$ 的多项式或有理式——高亏格振幅的多项式结构由此落地。
  \end{itemize}
\end{keytakeaways}

多项式构造中生成元的构造方式已在文献 \cite{ref26} 中讨论。出发点是选取模空间上的一个局部坐标 $z_*$，使
$C_{*ij}$ 为可逆的 $n\times n$ 矩阵，从而把式 (4.31) 改写为
\begin{equation}\tag{4.36}
S^{ij} = \left(C_*^{-1}\right)^{ik}
\left(\delta^j_* K_k + \delta^j_k K_* - \Gamma^j_{*k} + s^j_{*k}\right),
\end{equation}
式 (4.33) 中的自由度可用来选定部分 $s^k_{ij}$ \cite{ref26}。其余生成元则利用方程 (4.32) 构造 \cite{ref26}：
\begin{equation}\tag{4.37}
S^i = \frac{1}{2}\left(\partial_i S^{ii} - C_{imn}S^{mi}S^{ni}
     + 2s^i_{im}S^{mi} - h_i^{ii}\right),
\end{equation}
\begin{equation}\tag{4.38}
S = \frac{1}{2}\left(\partial_i S^i - C_{imn}S^m S^{ni}
   + s^i_{im}S^m + h_{im}S^{mi} - h_i^i\right).
\end{equation}
在这两个式子中指标 $i$ 是固定的，即对该指标不求和。式 (4.33) 中给生成元添加全纯函数的自由度，同样
可用来对函数 $h_i^{ii}$、$h_i^i$ 作某种选择；其余函数由这一选择固定，并可由式 (4.32) 算出。

% [Tutor Note: 物理图像] $S^{ij},S^i,S$ 就是 §4.3 费曼图图像中的传播子
\begin{yizhu}[物理图像]
生成元的物理意义直接来自第 4.3 节的费曼图图像：$S^{ij}$、$S^i$、$S$ 分别是 $K^{ij}=-S^{ij}$、
$K^{i\phi}=-S^i$、$K^{\phi\phi}=-2S$ 这些传播子，而 $C_{*ij}$ 可逆的局部坐标保证了可以“解出”传播子。
所谓“构造生成元”，就是选定一组传播子与全纯规范，使得建立反常方程所需的全部全纯函数都成为局部坐标
$z$ 的有理式乃至多项式——这正是高亏格振幅能写成多项式的技术前提。
\end{yizhu}

% [Tutor Note: 易错点] (4.36) 的固定指标 $*$ 与矩阵求逆的缩并方式
\begin{yizhu}[易错点]
式 (4.36) 的指标结构极易写错：$C_{*ij}$ 可逆，$(C_*^{-1})^{ik}$ 与 $C_{*kl}$ 在 $k$ 上缩并，故
$S^{ij}$ 的两个自由指标是 $i,j$；$\Gamma^j_{*k}$ 与 $s^j_{*k}$ 的固定指标 $*$ 不是求和指标。另一个
常见错误是把 $\delta^j_* K_k$ 与 $\delta^j_k K_*$ 互换——由 (4.31) 反解可知，前者的 $K$ 带下标 $k$
（与 $C$ 的第二指标缩并），后者带固定指标 $*$。建议用 $n=1$ 的最简情形逐项验证一次。
\end{yizhu}

% [Tutor Note: 易错点] (4.37)、(4.38) 中“指标 i 固定”的约定
\begin{yizhu}[易错点]
式 (4.37)、(4.38) 下方特意声明“指标 $i$ 固定，不求和”，这是全文最容易漏掉的约定：$S^{ii}$ 中的
$i$ 不求和，而 $C_{imn}S^{mi}S^{ni}$ 中 $m,n$ 求和。这两式是把闭合方程 (4.32) 沿某一固定方向（对
五次超曲面即唯一的 $z$ 方向）解出 $S^i$、$S$ 的结果，因此只能在该方向上有意义；若误按爱因斯坦约定
对 $i$ 求和，两式都会变成平凡恒等式而失去内容。
\end{yizhu}

\subsubsection{边界条件}
\label{sec:4.5.2}

\begin{keytakeaways}
  \begin{itemize}[leftmargin=1.4em,itemsep=2pt,topsep=2pt]
    \item \textbf{物理图像}：边界条件把 $\mathcal{F}^{(g)}$ 钉在模空间的边界上：大复结构极限处是
          常映射的拓扑贡献，类锥形流形轨迹处是额外无质量超多重态造成的奇异结构。
    \item \textbf{数学工具}：$g=1$ 的显式积分 (4.39)；常映射公式 (4.42)；锥形流形处的主导奇异与
          “间隙” (4.43)；全纯模糊性的局部 ansatz (4.44)。
    \item \textbf{关键结论}：$g=1$ 的系数 $s_i$、$r_a$ 被 (4.40)、(4.41) 完全固定；$g\ge2$ 时每个亏格
          只剩下 $p(z_i)$ 的有限个系数，由 (4.42)、(4.43) 给出的线性方程唯一确定。
  \end{itemize}
\end{keytakeaways}

% 原文此处为不编号的斜体小标题（\paragraph 级）
\medskip
\noindent\emph{亏格 1}\par

亏格 1 的全纯反常方程 (4.26) 可以积分给出
\begin{equation}\tag{4.39}
\mathcal{F}^{(1)}
= \frac{1}{2}\left(3+n-\frac{\chi}{12}\right)K
+ \frac{1}{2}\log\det G^{-1}
+ \sum_i s_i\log z_i
+ \sum_a r_a\log\Delta_a ,
\end{equation}
其中 $i=1,\ldots,n=h^{1,1}(Z)=h^{2,1}(Z^*)$，而 $a$ 遍历判别式的各个分支。系数 $s_i$ 与 $r_a$ 由
$\mathcal{F}^{(1)}$ 的主导奇异行为固定：对代数坐标 $z_i$，该行为由文献 \cite{ref19} 给出为
\begin{equation}\tag{4.40}
\mathcal{F}^{(1)} \sim -\frac{1}{24}\sum_i \log z_i \int_Z c_2 J_i ,
\end{equation}
% ==== 原文 PDF 第 36 页 ====
而对与锥形流形（conifold）奇点相对应的判别式 $\Delta$，主导行为由
\begin{equation}\tag{4.41}
\mathcal{F}^{(1)} \sim -\frac{1}{12}\log\Delta .
\end{equation}
给出。

% [Tutor Note: 易错点] (4.39) 的系数与 $\chi$ 的约定
\begin{yizhu}[易错点]
式 (4.39) 中 $K$ 的系数是 $\frac{1}{2}\left(3+n-\frac{\chi}{12}\right)$，注意这里的 $\chi$ 是靶空间
的欧拉示性数，$\chi=2\left(h^{1,1}-h^{2,1}\right)$，符号约定不同文献可能差一个整体符号；$n$ 同时
等于 $h^{1,1}(Z)$ 与 $h^{2,1}(Z^*)$。$\frac{1}{2}\log\det G^{-1}$ 项来自 $g=1$ 反常方程积分时的
Zamolodchikov 度规贡献，不要与 (4.42) 中的 $\frac{\chi}{2}$ 因子混淆。
\end{yizhu}

% [Tutor Note: 物理图像] (4.40) 与 (4.41) 如何固定 $s_i$、$r_a$
\begin{yizhu}[物理图像]
式 (4.40) 与 (4.41) 是固定 (4.39) 中待定系数的两条边界信息：前者给出 $z_i\to0$（大复结构/大体积）
处的对数发散，其系数由第二陈类 $c_2$ 沿曲线的积分给出；后者给出锥形流形轨迹 $\Delta=0$ 处的
$-\frac{1}{12}\log\Delta$ 发散。两者合起来给出了 $s_i$ 与 $r_a$ 的全部约束，因此 $g=1$ 的振幅可以被
完全确定——这正是 4.5.3 中 (4.46) 取值的来源。
\end{yizhu}

% 原文此处为不编号的斜体小标题（\paragraph 级）
\medskip
\noindent\emph{高亏格边界条件}\par

重构完整拓扑弦振幅所需的全纯模糊性，可以通过对 $\mathcal{F}^{(g)}$ 在模空间边界上施加各种边界条件
来固定。

% 原文此处为不编号的斜体小标题（\paragraph 级）
\medskip
\noindent\emph{大复结构极限}\par

$\mathcal{F}^{(g)}$ 在该点（它镜像于大体积极限）的主导行为已在文献 \cite{ref18,ref19,ref83,ref13,ref84,ref14} 中
算出。特别地，常映射的贡献为
\begin{equation}\tag{4.42}
\mathcal{F}^{(g)}\big|_{q_a=0}
= (-1)^g\frac{\chi}{2}\,
  \frac{|B_{2g}B_{2g-2}|}{2g(2g-2)(2g-2)!},
\qquad g>1,
\end{equation}
其中 $q_a$ 表示该点处指数化的镜像映射。

% [Tutor Note: 物理图像] 常映射贡献是 $F^{(g)}$ 中唯一与模无关的部分
\begin{yizhu}[物理图像]
式 (4.42) 是常映射的贡献：世界面不携带任何手征插入、整体映射到靶空间一点。它只依赖靶空间的
拓扑（$\chi$）与亏格，是固定亏格下唯一与模参数无关的项，因此给出 $\mathcal{F}^{(g)}$ 在大复结构
极限处的一个“锚点”。$B_{2g}$ 与 $B_{2g-2}$ 的出现来自常数映射下对亏格 $g$ 世界面模空间的积分
（所有世界面都映到同一点，只剩下模空间上的特征类积分），其组合
$|B_{2g}B_{2g-2}|/\left(2g(2g-2)(2g-2)!\right)$ 也是 $g$ 很大时 $\mathcal{F}^{(g)}$ 增长的主要来源，
这使该边界条件对高亏格递推非常强。
\end{yizhu}

% 原文此处为不编号的斜体小标题（\paragraph 级）
\medskip
\noindent\emph{类锥形流形轨迹}\par

配分函数 $\mathcal{F}^{(g)}$ 在锥形流形轨迹处的主导奇异行为已在文献 \cite{ref18,ref19,ref85,ref86,ref13,ref14} 中
确定：
\begin{equation}\tag{4.43}
\mathcal{F}^{(g)}(t_c)
= b\,\frac{B_{2g}}{2g(2g-2)\,t_c^{2g-2}} + O(t_c^0),
\qquad g>1 .
\end{equation}
这里 $t_c\sim\Delta^{1/m}$ 是判别式轨迹 $\Delta=0$ 处的平坦坐标。对于锥形流形奇点，$b=1$、$m=1$。
特别地，(4.43) 中的主导奇异性以及次主导奇异项的缺失，来自文献 \cite{ref13,ref14} 的 Schwinger 圈计算，该计算
给出时空理论中额外无质量超多重态（hypermultiplet）的效应 \cite{ref87}。这种奇异结构以及次主导奇异项的
“间隙”（gap），也在对偶矩阵模型中观察到 \cite{ref88}，并首先被文献 \cite{ref89,ref90} 用来固定高亏格的全纯模糊性。
文献 \cite{ref13,ref14} 的时空推导并不限于锥形流形情形，也适用于 $m>1$ 的奇点，后者在时空中给出由额外无质量
矢量多重态与超多重态组成的不同谱。Schwinger 圈积分的系数是对粒子自旋的加权求迹 \cite{ref87,ref86}，由此预言
主导奇异项的系数为 $b=n_H-n_V$。前因子（prefactor）$b$ 在 $m>1$ 情形下的出现，已在文献 \cite{ref26} 中就
局部 $\mathbb{F}_2$ 的情形讨论过（另见 \cite{ref91}）。

% [Tutor Note: 易错点] (4.43) 的“间隙”性质给出多个线性条件
\begin{yizhu}[易错点]
式 (4.43) 有两个易错点。其一，$t_c\sim\Delta^{1/m}$ 中的 $m$ 不一定为 $1$；$m>1$ 时前因子为
$b=n_H-n_V$，需要用局部模型的谱重新计算，而不是简单取 $b=1$。其二，“次主导奇异项的缺失”是指
在主导项 $t_c^{2-2g}$ 之后，所有更奇异的项（如 $t_c^{1-2g}$ 等）都缺席（gap）；$O(t_c^0)$ 表示允许
常数项与正规项，并非“全部为零”。正因为有这个间隙，一次边界条件能同时钉死 $p(z_i)$ 中的一批系数；
把它误读成“所有更高项为零”会导致系数超定、方程组无解。
\end{yizhu}

% ==== 原文 PDF 第 37 页 ====

% 原文此处为不编号的斜体小标题（\paragraph 级）
\medskip
\noindent\emph{全纯模糊性}\par

$\mathcal{F}^{(g)}$ 的奇异行为通过全纯模糊性的如下局部 ansatz 加以考虑：
\begin{equation}\tag{4.44}
\mathrm{hol.ambiguity} \sim \frac{p(\tilde z_i)}{\Delta^{(2g-2)}} ,
\end{equation}
其中 $p(\tilde z_i)$ 先验地是奇点附近局部坐标 $\tilde z_i$ 的级数。然而，把各奇点处的局部信息与
具有有限单值性的边界除子拼接起来，可知分子 $p(z_i)$ 一般只是 $z_i$ 的低次多项式。这里 $z_i$ 表示
以大复结构为中心的自然坐标，$z_i=0\ \forall i$。$p(z_i)$ 中有限个系数受 (4.43) 约束。高亏格振幅的
多项式结构结合这些边界条件，足以在局部 CY 流形上微扰地求解拓扑弦理论（例如见文献 \cite{ref92,ref26}），也使
得在紧致 CY 流形上求解高亏格振幅成为可能 \cite{ref90,ref91,ref28}。

% [Tutor Note: 含糊点] (4.44) 中 $\tilde z_i$ 与 $z_i$ 的记号切换
\begin{yizhu}[含糊点]
式 (4.44) 及其上一句用局部坐标 $\tilde z_i$ 写 ansatz，紧接着的讨论却改用大复结构坐标 $z_i$，原文
未明说这是同一个对象的两种坐标。可这样理解：$p$ 先在奇点附近的局部坐标中是级数，拼接（patching）
到全局并把各边界除子的信息纳入后，其分子在自然坐标 $z_i$ 下成为低次多项式；$z_i$ 正是 (4.47) 中
$P=\Delta^{-1}$ 所用的坐标。核对 (4.44) 与 (4.47)、(4.48) 时，须把 $p$ 理解为 $z$ 的多项式。
\end{yizhu}

\subsubsection{例：五次超曲面}
\label{sec:4.5.3}

\begin{keytakeaways}
  \begin{itemize}[leftmargin=1.4em,itemsep=2pt,topsep=2pt]
    \item \textbf{物理图像}：五次超曲面是最简单的单参数例子（$h^{1,1}=1$、$\chi=-200$），用来演示
          多项式结构如何变成逐亏格的显式计算。
    \item \textbf{数学工具}：全纯函数的指标乘除 $z$ 约定、按文献 \cite{ref25} 选取的 $s$、$h$ 值 (4.45)、
          亏格 1 的边界数据 (4.46)、以及由此固定的 $z\partial_z\mathcal{F}^{(1)}$ (4.47)。
    \item \textbf{关键结论}：$g=1$ 与 $g=2$ 的振幅都被完全固定，其中 $g=2$ 的显式结果 (4.48) 是多项式
          结构可用性的直接检验；用同一套办法可以走到很高亏格 \cite{ref90}。
    \item \textbf{可核对性}：$P=\Delta^{-1}$ 视为 $0$ 次，则 (4.48) 中每一项的总次数都不超过
          $3g-3=3$，这是检验抄录与理解是否正确的快捷判据。
  \end{itemize}
\end{keytakeaways}

% 原文此处为不编号的斜体小标题（\paragraph 级）
\medskip
\noindent\emph{生成元的选取}\par

多项式生成元的选取方式是固定的：使式 (4.31)、(4.32) 中出现的全纯函数都是 $z$ 的有理表达式，其中
$z$ 是模空间大复结构片中的坐标。对于这些全纯函数，下列上下指标分别乘以、除以 $z$：
\[
  A^j_i \to \frac{z_i}{z_j}A^j_i .
\]
采用这一约定后，多项式构造框架中出现的所有全纯函数都可以用局部坐标的多项式表示（见文献
\cite{ref21,ref22,ref25}）。这些全纯函数按文献 \cite{ref25} 选取为
\begin{equation}\tag{4.45}
s^z_{zz} = -\frac{8}{5},
\qquad
h^{zz}_z = \frac{1}{25},
\qquad
h^z_z = -\frac{1}{125},
\qquad
h_z = \frac{2}{3125},
\qquad
h_{zz} = \frac{2}{25},
\end{equation}

% [Tutor Note: 易错点] 指标乘除 $z$ 的约定与 (4.45) 的归一化
\begin{yizhu}[易错点]
$A^j_i\to(z_i/z_j)A^j_i$ 这一约定不是装饰：它把不同权重的全纯函数按 $z$ 的幂次归一，使得 (4.45) 中
的常数成为纯数字。若漏掉该约定，或改用另一套 $z$ 的归一化（例如把大复结构坐标乘以常数），(4.45) 的
数值就会改变。核对时请同时对照文献 \cite{ref25} 的坐标定义与 (4.47) 中 $P=\Delta^{-1}$ 的定义；这三处必须
用同一个 $z$。
\end{yizhu}

% 原文此处为不编号的斜体小标题（\paragraph 级）
\medskip
\noindent\emph{亏格 1}\par

利用亏格 1 的边界条件，可以如下固定关于 $\mathcal{F}^1$ 的方程 (4.39)：
\begin{equation}\tag{4.46}
n = 1, \qquad \chi = -200, \qquad s = -\frac{31}{12}, \qquad r = -\frac{1}{12},
\end{equation}
而起始的多项式表达式被固定为（其中 $P=\Delta^{-1}$）：
\begin{equation}\tag{4.47}
z\partial_z\mathcal{F}^{(1)}
= \frac{1}{60}\left(-112 + 560K_z + 5P(1+30S^{zz})\right).
\end{equation}

% [Tutor Note: 物理图像] (4.46) 的 $\chi=-200$ 与五次超曲面的霍奇数
\begin{yizhu}[物理图像]
式 (4.46) 的数值都可以独立核对：五次超曲面有 $h^{1,1}=1$、$h^{2,1}=101$，故
$\chi = 2(h^{1,1}-h^{2,1}) = 2(1-101) = -200$，$n=h^{1,1}=1$；而 $s=-\frac{31}{12}$、
$r=-\frac{1}{12}$ 分别由 (4.40) 在 $z\to0$ 的主导行为与 (4.41) 在 $\Delta\to0$ 的主导行为确定。
也就是说，在 (4.39) 的形式下 $g=1$ 的振幅不再含由 $s_i$、$r_a$ 表示的自由参数——它们已被边界条件
完全固定。
\end{yizhu}

% [Tutor Note: 易错点] (4.47) 中 $P$ 不是带分次的生成元
\begin{yizhu}[易错点]
式 (4.47) 与 (4.48) 中的 $P=\Delta^{-1}$ 是 $z$ 的全纯函数，不是 $K_i,S^{ij},S^i,S$ 那类带分次的
生成元；在次数核算中应把它按 $0$ 次处理。这解释了为何 (4.48) 中 $P^2$ 乘以三次项后总次数仍为
$3g-3=3$。另一个易错点是 (4.47) 左边是 $z\partial_z\mathcal{F}^{(1)}$ 而非 $\mathcal{F}^{(1)}$ 本身，
这是为了配合大复结构坐标下的多项式表示；抄录时注意 $P$ 项是 $5P(1+30S^{zz})$，括号内的 $1$ 不要
漏掉。
\end{yizhu}

% ==== 原文 PDF 第 38 页 ====

% 原文此处为不编号的斜体小标题（\paragraph 级）
\medskip
\noindent\emph{高亏格}\par

利用多项式形式 (4.34)、(4.35) 的全纯反常方程以及边界条件，可以固定高亏格振幅。例如在亏格 2、取上述
生成元时，得到
\begin{equation}\tag{4.48}
\begin{aligned}
\mathcal{F}^{(2)}
&= \frac{1}{36000}\Bigl(
  5P^2\left(37500(S^{zz})^3 + 6000(S^{zz})^2 + 325S^{zz} + 6\right) \\
&\qquad - P\left(25000S^z(30S^{zz}+1) + 150000(S^{zz})^2 + 4400S^{zz} - 23\right) \\
&\qquad + 16\left(175000S + 35000S^z + 3500S^{zz} - 39\right)\Bigr).
\end{aligned}
\end{equation}

% [Tutor Note: 物理图像] 用 Ref.\cite{ref21} 生成元加边界条件把五次解到高亏格的逻辑
\begin{yizhu}[物理图像]
(4.48) 是整套算法的样板：先用上面构造的生成元把 $\mathcal{F}^{(2)}$ 的非全纯部分由反常方程定出，
剩下的全纯模糊性按 (4.44) 型 ansatz 只有有限个系数；再用 (4.42) 的常映射项与 (4.43) 的间隙条件逐项
比较，恰好得到与未知系数个数相同的线性方程，解出的就是 (4.48) 中的那些数字（$36000$、$37500$、
$6000$、$325$、$25000$、$150000$、$4400$、$175000$、$35000$、$3500$、$39$ 等）。文献 \cite{ref90} 正是用
文献 \cite{ref21} 的多项式生成元、前面综述的边界条件，再加文献 \cite{ref93} 的不变量计数技巧，把五次超曲面解到了
很高的亏格。
\end{yizhu}

% [Tutor Note: 易错点] (4.48) 的逐项核对清单
\begin{yizhu}[易错点]
抄录 (4.48) 时建议按结构而非按行核对：整体是 $\frac{1}{36000}\left(5P^2A - PB + 16C\right)$，其中
$A=37500(S^{zz})^3+6000(S^{zz})^2+325S^{zz}+6$（三次到零次），
$B=25000S^z(30S^{zz}+1)+150000(S^{zz})^2+4400S^{zz}-23$（三次到零次），
$C=175000S+35000S^z+3500S^{zz}-39$（三次到零次）。三处最易出错：$B$ 中 $30S^{zz}+1$ 的括号范围、
$B$ 末项是 $-23$ 而非 $+23$、$C$ 末项是 $-39$；此外 $16$ 乘在整个 $C$ 上，不要漏乘。
\end{yizhu}

% ---------------------------------------------------------------------------
\section*{致谢}
\addcontentsline{toc}{section}{致谢}
% ---------------------------------------------------------------------------

我要感谢清华大学数学科学中心（Mathematical Sciences Center, Tsinghua University）与
Shing-Tung Yau 教授的盛情接待，以及邀请我作这些讲座。我还要感谢 Michael Hecht、Albrecht Klemm、
Dominique Länge、Peter Mayr、Emanuel Scheidegger、Shing-Tung Yau 与 Jie Zhou，就本讲义所涵盖的
某些课题及相关问题所作的讨论与合作。此外，我要感谢 Emanuel Scheidegger 与 Jie Zhou 对书稿的评论。
本工作由 DFG fellowship AL 1407/1-1 资助。

% [Tutor Note: 含糊点] 致谢中的外文姓名与资助编号照录
\begin{yizhu}[含糊点]
致谢中的人名按原文保留拉丁字母。Dominique Länge 的 “ä” 是德语变音，用 XeLaTeX 直接以 UTF-8 写入
即可，不要写成 “Lange”。资助编号 “AL 1407/1-1” 是德国研究基金会（DFG）的 fellowship 编号，属
专名，照抄不译。另外，原文的 “Prof. Shing-Tung Yau” 在两句感谢中重复出现，译文照原文保留这一重复。
\end{yizhu}

% ---------------------------------------------------------------------------
% 术语建议（供 Lead 统一时参考，未在正文中混用）：
%   automorphic forms = 自守形式；almost holomorphic modular forms = 近全纯模形式；
%   grading = 分次；hol.ambiguity = 全纯模糊性（原文用缩写，正文照录）。
% 版式说明：4.5.2、4.5.3、致谢中的不编号斜体小标题（Genus 1、The large complex
%   structure limit 等）原文为 \paragraph 级斜体标题，本片段用“斜体独立行”复现，
%   未新增编号层级；若全书统一改用 \paragraph* 或 \subsubsection*，可直接替换。
% ---------------------------------------------------------------------------
